% !TeX root = ../../main.tex
\section{选择性执行的成本结构与稳健性讨论}
\subsection{簇代表成本与成员扫描成本}
选择性检索增加了簇代表比较。设簇数为 $C$，代表矩阵非零数为 $z_c$，选中成员的非零数为 $z_s$，则在线稀疏工作量至少包含 $z_c+z_s$。如果簇代表趋于稠密或簇数过多，即使成员扫描量下降，代表计算也可能成为不可忽略的固定成本。

采用实际稀疏成员作为代表能够约束这一成本，但可能牺牲均值代表的拟合能力。选择代表时需联合观察代表非零数、簇内方向凝聚度及达到目标召回的探测数。仅比较聚类目标函数不足以确定检索系统性能，因为代表稀疏性与在线执行成本并不直接出现在普通 K-means 损失中。

\begin{table}[htbp]\centering\caption{候选簇粒度对系统成本的双向影响}
\begin{tabularx}{\textwidth}{p{2.6cm}XX}\toprule
变化&潜在收益&潜在代价\\\midrule
簇数量增加&更精细的范围选择&代表比较和地址表增长\\
簇数量减少&减少索引与探测固定开销&单簇扫描范围变大\\
代表更稀疏&降低代表访存&可能降低方向代表性\\
簇内更紧凑&相关候选更集中&离线细化成本增加\\
簇大小更均匀&降低核间极端负载&可能偏离相似度最优划分\\\bottomrule
\end{tabularx}\end{table}

\subsection{簇级召回与元素剪枝的交互}
簇级选择排除整组候选，元素剪枝改变保留候选的近似得分。两种误差作用位置不同：被排除簇中的真实 Top-K 无法通过提高簇内精度恢复，而被量化扰乱的候选可能通过更多预算或原始得分重算改善。因此，调参应先确认簇范围有足够覆盖，再讨论簇内表示精度。

设所探测簇包含的真实 Top-K 比例为 $R_c$，在这些真实候选中经过局部筛选仍保留的条件比例为 $R_l$，则完整候选保持可写为 $R_cR_l$。这里是基于同一查询集合事件划分的集合计数关系，不是假设两种误差独立。对跨查询平均值，平均乘积通常不等于两个平均值的乘积，故不应通过分别测得的平均召回直接相乘估算端到端结果。

这个分解可以指导消融：一组实验以原始得分在选中簇内精确排序，测量范围选择的上限；另一组固定同一候选范围，比较 INT6 和 ReSparse 后的结果。最后再启用分布式候选维护，检查硬件近似是否增加损失。完整结果仍以端到端测量为准。

\subsection{探测预算取整与核内选择}
32 核系统将全局预算分解到本地执行。若 $n_{\mathrm{probe}}$不是32 的倍数，统一向上取整会使实际预算最多多出31 个簇；若某核簇数不足，还需截断其预算。因此，命令参数并不必然等于真正访问簇数。实验应记录实际探测簇、实际候选数与实际非零数三个量。

\begin{table}[htbp]\centering\caption{探测预算的三个层次}
\begin{tabularx}{\textwidth}{p{2.5cm}XX}\toprule
指标&含义&为何不能互相替代\\\midrule
请求探测数&软件传入的预算&受到分布式分解和取整影响\\
实际簇数&所有核真正访问簇的数量&不同簇包含的成员不同\\
实际行数&参与相似度计算的向量数量&不同向量非零数不同\\
实际非零数&需要传输和乘加的记录数量&还未包含索引和事务开销\\\bottomrule
\end{tabularx}\end{table}

局部预算也可能造成相关簇集中区域被限制。增加每核预算能够缓解，但会在不相关核上增加访问。更灵活的预算分配可以基于代表得分或历史需求调整，然而需要额外协调信息，与完全独立 ACPE 执行存在取舍。本文采用固定分解获得较简单控制路径，并通过实测召回评价其结果。

\subsection{热点画像的时间尺度}
热点画像是对查询分布的统计估计。画像查询过少，容易受偶然候选影响；过长的历史窗口则可能掩盖近期变化。对于静态数据但查询意图变化的服务，簇结构仍有效而热点映射可能失效，因此索引重建与映射更新应视为不同层次的维护操作。

设新映射预计每次查询节省 $\Delta t$，迁移成本为 $T_m$，预测在下一次明显分布变化前还能处理 $N$个查询，则需要 $N\Delta t>T_m$才能获得总时间收益。这里的节省还应包含迁移期间的干扰，而不是只比较迁移前后的稳定状态 QPS。

原论文的受控正偏置实验验证一种特定分布变化，不代表所有真实漂移。后续可区分查询主题变化、范数变化、稀疏度变化与突发热点，分别检查 CV、最大核时间和尾延迟。只有把漂移类型与质量保持一起记录，才能判断映射是否真正稳健。

\subsection{长尾簇与不可分割任务}
贪心分配能均衡多个簇的累计工作，但若一个簇本身的工作量远超平均值，即使独占一个核也会形成瓶颈。将任务视为不可分割单元时，理论完成时间至少为最大单簇负载与平均负载两者的最大值。这一界解释了仅靠重新放置无法解决所有热点。

拆分热点簇能够增加调度自由度，却可能改变代表集合和探测语义；复制热点数据则需要同步更新并消耗容量。两者都需要重新定义结果合并和去重规则。如果复制后同一候选可由多个核返回，应在最终 Top-K 前去重，防止重复 ID 挤占结果容量。

\subsection{性能曲线与前沿比较方法}
对每个系统，通过多个探测预算获得 $(R,Q)$点，其中 $R$为召回、$Q$为 QPS。在相同或更高召回下仍具有更高吞吐的点，支配被比较点。可移除被支配点构成观测前沿，但这只覆盖实际测量的参数空间，不等价于全部配置的最优边界。

如果目标召回位于两个测量点之间，可以报告达到或超过目标的最小预算点，或明确说明使用插值。不同系统若使用不同插值方法，可能夸大差异。对固定全量系统，其质量和吞吐是固定点或近似水平参考，不应将其误画成能够连续调节召回的算法。

还应区分算术平均 QPS、几何平均加速比和最大加速比。前者受绝对吞吐规模影响，后者衡量相对比值，最大值只描述最有利案例。本文分别报告上述统计量，并明确其对应的数据集合与质量区间。

\subsection{尾延迟、能耗与在线服务}
平均 QPS 无法直接描述尾延迟。同步多核查询由最晚完成的核决定，偶发热点会使少量查询耗时较长；系统排队又可能进一步放大尾部。在离线单查询测量中观察到的 P99 改善，不自动等同于固定到达率服务下的 P99 改善。

在线服务评估需要给出到达过程、并发数和预热方式，并区分设备执行时间与排队时间。能耗应对整个测量窗口积分，而不是仅使用标称功耗乘以一个最短内核时间。若 CPU 负责候选合并或索引维护，其成本也应纳入相应的系统边界。

\subsection{高维数据的适配边界}
当前 ACPE 继承10 位列索引和多份片上查询缓存，适合论文中的维度范围。对数万维稀疏词项向量，列索引变宽会减少每包非零数量，完整复制查询也可能不可行。可考虑分块查询缓存或稀疏查询索引，但会引入额外间接访问和缓存未命中。

因此，将本系统拓展到高维词项检索需重新评估数据格式、簇代表非零度和查询缓存路径，而不是只提高一个维度参数。该边界也说明实验材料中的经典词向量稀疏化数据与现代神经稀疏检索数据应分别取得，保持论证范围清楚。
